% Part: turing-machines
% Chapter: undecidability
% Section: introduction

\documentclass[../../../include/open-logic-section]{subfiles}

\begin{document}

\olfileid{tur}{und}{int}
\olsection{مقدمة}

قد يبدو بديهيًا ألا تكون كل دالة، بل حتى كل دالة حسابية، قابلة
للحساب. فعدد الدوال هائل وسلوكها بالغ التعقيد؛ ومن أمثلتها الدوال
المعرّفة انطلاقًا من اضمحلال الجسيمات المشعة أو من غيره من السلوك
الفوضوي أو العشوائي. لنفترض أننا نبدأ عد فترات طول كل منها ثانية واحدة
من لحظة معينة، ونعرّف الدالة $f(n)$ بأنها عدد الجسيمات في الكون التي
تضمحل في الفترة ذات الثانية الواحدة رقم~$n$ بعد تلك اللحظة الابتدائية.
تبدو هذه الدالة مرشحة لأن تكون دالة لا أمل لنا قط في حسابها.

لكن العجز عن تصور كيفية حساب مثل هذه الدوال شيء، والبرهنة بالفعل على
أنها غير قابلة للحساب شيء آخر تمامًا. بل قد تكون الدوال التي تبدو
بالغة التعقيد إلى حد اليأس قابلة للحساب بمعنى مجرد. لنفترض، مثلًا، أن
عمر الكون منتهٍ، وأنه سينكمش يومًا في المستقبل البعيد جدًا إلى نقطة
واحدة، كما تتنبأ بعض النظريات الكونية. عندئذ لا يوجد إلا عدد منتهٍ
(وإن كان هائلًا على نحو لا يصدق) من الثواني اللاحقة لتلك اللحظة
الابتدائية التي تكون $f(n)$ معرّفة عندها. وكل دالة لا تكون معرّفة إلا
عند عدد منتهٍ من المدخلات قابلة للحساب: إذ يمكننا إدراج مخرجاتها في
جدول واحد كبير، أو ترميزها في مخطط انتقال حالات بالغ الكبر لآلة تورنغ.

نهتم كثيرًا بحالات خاصة من الدوال تعطي قيمها أجوبة عن أسئلة نعم/لا.
فمثلًا يقترن السؤال «هل $n$ عدد أولي؟» بالدالة
\[
\fn{isprime}(n) = \begin{cases}
  1 & \text{إذا كان $n$ أوليًا}\\
  0 & \text{خلاف ذلك.}
  \end{cases}
\]
نقول إن سؤال نعم/لا \emph{قابل للتقرير فعليًا} إذا كانت الدالة
المقترنة به، ذات القيمتين $1/0$، قابلة للحساب فعليًا.

لكي نبرهن رياضيًا أن ثمة دوال لا يمكن حسابها فعليًا، أو مسائل لا يمكن
تقريرها فعليًا، لا بد من تثبيت نموذج معين للحساب، ثم بيان أن ثمة دوال
لا يستطيع حسابها أو مسائل لا يستطيع تقريرها. ويمكننا، مثلًا، أن نبين
أن آلات تورنغ لا تستطيع حساب كل دالة، ولا تستطيع تقرير كل مسألة. ثم
يمكننا الاحتكام إلى أطروحة تشرتش--تورنغ لنستنتج أن آلات تورنغ ليست
وحدها عاجزة عن حساب كل دالة، بل لا يستطيع ذلك أي إجراء فعلي.

مفتاح البرهنة على مثل هذه النتائج السلبية هو إمكان إسناد أعداد إلى
آلات تورنغ نفسها. وأسهل طريقة لذلك هي تعدادها، ربما بتثبيت طريقة معينة
لكتابة آلات تورنغ وبرامجها، ثم إدراجها إدراجًا منهجيًا. ومتى رأينا أن
هذا ممكن، نتج وجود دوال غير قابلة للحساب بآلة تورنغ من اعتبارات بسيطة
تتعلق بالأعداد الأصلية: فمجموعة الدوال من $\Nat$ إلى~$\Nat$ (بل حتى
مجرد الدوال من $\Nat$ إلى $\{0, 1\}$) !!{nonenumerable}، لكن مجموعة
الدوال القابلة للحساب بآلة تورنغ ليست إلا !!{denumerable}، لأننا نستطيع
تعداد جميع آلات تورنغ.

ويمكننا أيضًا تعريف دوال ومسائل \emph{معينة} نستطيع البرهنة على أنها
غير قابلة للحساب وغير قابلة للتقرير، على الترتيب. ومن هذه المسائل ما
يسمى \emph{مسألة التوقف}. يمكن وصف آلات تورنغ وصفًا منتهيًا بإدراج
تعليماتها. ويمكن بالطبع استعمال مثل هذا الوصف لآلة تورنغ، أي برنامج
آلة تورنغ، دخلًا لآلة تورنغ أخرى. ومن ثم نستطيع النظر في آلات تورنغ
تقرر أسئلة عن آلات تورنغ أخرى. ومن الأسئلة المثيرة للاهتمام بوجه خاص:
«هل تتوقف آلة تورنغ المعطاة في نهاية المطاف حين تبدأ عند الدخل~$n$؟»
وسيكون جميلًا لو وجدت آلة تورنغ تستطيع تقرير هذا السؤال: تخيلها آلة
تورنغ لضبط الجودة تضمن ألا تعلق آلات تورنغ في حلقات غير منتهية وما
شابه. لكن الحقيقة المثيرة التي برهن عليها تورنغ هي استحالة وجود مثل
هذه الآلة. فلا يمكن أن توجد آلة تورنغ واحدة تبدأ بدخل مؤلف من وصف آلة
تورنغ~$M$ وعدد~$n$، ثم تتوقف دائمًا بخرج $1$ أو $0$ بحسب ما إذا كانت
$M$ ستتوقف عند بدء تشغيلها على الدخل~$n$ أم لا.

ومتى كان لدينا أمثلة لمسائل معينة غير قابلة للتقرير، استطعنا استعمالها
لبيان أن مسائل أخرى غير قابلة للتقرير أيضًا. ومن المسائل الشهيرة غير
القابلة للتقرير، مثلًا، السؤال: «هل صيغة الرتبة الأولى~$!A$ صادقة
منطقيًا؟». فلا توجد آلة تورنغ مضمونة التوقف، عند إعطائها صيغة رتبة
أولى~$!A$ دخلًا، بخرج $1$ أو $0$ بحسب ما إذا كانت $!A$ صادقة منطقيًا
أم لا. وكان سؤال إيجاد إجراء يحل هذه المسألة فعليًا يسمى تاريخيًا
«مسألة القرار» بإطلاق؛ ولذلك نقول إن مسألة القرار غير قابلة للحل. وقد
برهن تورنغ وتشرتش هذه النتيجة كل على حدة في الوقت نفسه تقريبًا، ولذلك
تسمى أيضًا مبرهنة تشرتش--تورنغ.

\end{document}
